\chapter{Resultados y limitaciones}
\section{Artefactos preparados}
La sesión construye el primer incremento local y el conjunto documental: diagnóstico, requisitos, backlog de 28 historias, criterios de aceptación, trazabilidad, proceso, planificación, decisiones, diagramas y fuentes del informe. El alcance funcional de implementación se concentra en autenticación, usuarios, roles y capacitaciones.

La evidencia de ejecución se lee en la tabla~\ref{tab:evidencia}. Un archivo creado, una ruta escrita o un botón visible no prueban por sí solos que el flujo esté verificado. Los resultados finales deben indicar los comandos efectivamente completados y sus limitaciones.

\section{Resultado esperado del flujo V1}
El administrador de demostración puede trabajar con cuentas ficticias y capacitaciones almacenadas en SQLite a través de una API protegida. La interfaz presenta carga, confirmación y error según la respuesta real. El acceso y la edición dependen de los permisos actuales, no de un rol escrito manualmente en el navegador.

Se preservan las ramas de integrantes y los documentos previos. No se presenta la construcción local como una integración, publicación, revisión de compañeros o aprobación de Paz Sur. No se crean tags V1--V4 para simular una evolución que no ocurrió.

\section{Limitaciones del producto}
\begin{itemize}
  \item No están implementados los flujos completos de participantes, asistencia, sesiones, resultados o historial.
  \item No se emiten certificados, ni se presupone que todas las capacitaciones los necesitan.
  \item Los catálogos de marcas, campañas, sedes y áreas son diseño candidato.
  \item Los reportes consolidados y las alertas no se presentan como disponibles por existir un permiso de catálogo.
  \item La matriz de roles, estados y reglas académicas necesita validación con el cliente.
  \item No se ha demostrado capacidad de producción, alta disponibilidad o cumplimiento normativo específico.
\end{itemize}

\section{Limitaciones del entorno y del proceso}
La falta de espacio en la unidad del perfil complicó la ejecución de herramientas. Se verificó disponibilidad de Node y SQLite. XeLaTeX se ejecutó realmente con instalación automática deshabilitada y falló por ausencia de \texttt{xelatex.fmt}; no pudo completar la generación del formato. Se preparó la alternativa pdfLaTeX con fuentes Latin Modern vectoriales y se corrigió la interacción entre Babel español y TikZ. El resultado de esa compilación se consigna en la tabla de evidencia, separado de la revisión visual del PDF.

Los Mermaid se revisaron como fuentes, pero su renderizado con un motor Mermaid no está verificado. El informe incluye esas fuentes y dos diagramas nativos TikZ; ambos tipos de artefacto deben distinguirse al presentar resultados.

Las ceremonias Scrum, las revisiones de compañeros y la integración en Git permanecen propuestas o pendientes hasta que existan registros. Las estimaciones de capacidad son propuestas y no reemplazan datos históricos de Jira.

\section{Siguiente actividad recomendada}
El equipo debe revisar el diff y ejecutar la demostración con datos ficticios, validar las reglas del primer incremento y registrar las observaciones. Después podrá seleccionar una tarea pequeña de participantes o instructores, verificar sus dependencias y completar sus criterios sin comprometer automáticamente el resto del producto.
